% Options for packages loaded elsewhere
\PassOptionsToPackage{unicode}{hyperref}
\PassOptionsToPackage{hyphens}{url}
%
\documentclass[11pt]{article}
\usepackage[top=0.9in, bottom=0.9in, left=0.9in, right=0.9in]{geometry}
\usepackage{enumitem}
\usepackage{titlesec}
\usepackage{fancyvrb}
\usepackage{etoolbox}
\setlist[enumerate]{itemsep=2pt, topsep=4pt, parsep=0pt}
\setlist[itemize]{itemsep=2pt, topsep=4pt, parsep=0pt}
\setlist[description]{itemsep=2pt, topsep=4pt, parsep=0pt}
\setlength{\parskip}{0pt}
\setlength{\parindent}{0pt}
\renewcommand{\baselinestretch}{0.95}
\titlespacing*{\section}{0pt}{6pt}{3pt}
\titlespacing*{\subsection}{0pt}{5pt}{2pt}
\titlespacing*{\subsubsection}{0pt}{4pt}{2pt}
\let\oldverbatim=\verbatim
\let\oldendverbatim=\endverbatim
\renewenvironment{verbatim}{\vspace{-6pt}\oldverbatim}{\oldendverbatim\vspace{-4pt}}
\begin{document}
\maketitle

\section*{Abstract}
DataFrameLib is a C++17 DataFrame library built on Apache Arrow. It provides
two execution modes: eager (Pandas-like, immediate) and lazy (Polars-like, 
optimized). A fixed-point query optimizer implements five standard SQL rewrites: 
predicate pushdown, projection pushdown, constant folding, expression 
simplification, and limit pushdown. The codebase is compact (~2500 LOC) but 
complete: it demonstrates a working eager API, a logical plan representation, 
an optimizer, and an executor that reuses eager operators for lazy queries.

\section{Overall Design}

DataFrameLib is structured as three cooperating layers (Figure 1):

\begin{enumerate}
  \item \textbf{EagerDataFrame}: immediate execution. A thin, immutable wrapper 
    around \texttt{std::shared\_ptr<arrow::Table>}. Operations call Arrow Compute 
    kernels and return instantly.
  \item \textbf{LazyDataFrame}: deferred execution. Builds a logical plan DAG 
    of \texttt{LogicalNode} objects. Operations extend the DAG and return a new 
    lazy frame. Nothing executes until \texttt{collect()} is called.
  \item \textbf{QueryOptimizer}: optimizes the lazy plan. Applies five rewrite 
    rules to a fixed point, then hands the plan to a DFS executor that 
    materializes the result via eager operators.
\end{enumerate}

The expression system is a small AST rooted at \texttt{Expr}: \texttt{ColExpr}, 
\texttt{LitExpr}, \texttt{BinaryExpr}, \texttt{UnaryExpr}, \texttt{AggExpr}, 
\texttt{StringExpr}, \texttt{AliasExpr}. User code interacts via the 
operator-overloaded \texttt{ExprBuilder} and helper functions \texttt{col(name)} 
and \texttt{lit(value)}. Both eager and lazy modes share the same evaluator, 
which is key: it enables the optimizer to constant-fold expressions by 
evaluating them against a synthetic 1-row table.

\section{Arrow Integration}

Arrow is the physical backbone. Three properties drove the choice:

\begin{enumerate}
  \item \textbf{Columnar layout.} Each column is an 
    \texttt{arrow::ChunkedArray}, a sequence of contiguous primitive or 
    Utf8 buffers. Vectorized compute (sum, equals, greater-than, etc.) 
    benefits from cache locality and auto-vectorization in Arrow kernels.
  \item \textbf{Zero-copy sharing.} \texttt{arrow::Table} is immutable and 
    reference-counted. Our \texttt{EagerDataFrame} copy is an atomic pointer 
    copy. Operations like \texttt{head}, \texttt{select}, and \texttt{filter} 
    produce new tables but share unchanged buffers.
  \item \textbf{Type safety.} Arrow's \texttt{DataType} hierarchy is mapped 
    1:1 onto our \texttt{ColType} enum via \texttt{arrowTypeToColType}. 
    Unsupported Arrow types (e.g., \texttt{decimal128}) are rejected at load 
    time; we never silently coerce.
\end{enumerate}

Null handling is inherited verbatim. Every Arrow kernel propagates nulls in SQL 
semantics: \texttt{NULL + 1 = NULL}, \texttt{NULL > 5 = NULL} (dropped by 
filter), etc. Composing kernels gives null-safety ``for free''.

I/O uses \texttt{arrow::csv::TableReader} and 
\texttt{parquet::arrow::*}. File handles are RAII-wrapped to ensure automatic 
cleanup on scope exit.

\section{Expression System and Type Checking}

The expression AST is minimal but expressive:

\begin{description}
  \item[ColExpr] Column reference by name.
  \item[LitExpr] Scalar literal (int, float, string, bool).
  \item[BinaryExpr] Binary op (arithmetic, comparison, logical).
  \item[UnaryExpr] Unary op (negate, NOT, abs, is\_null).
  \item[StringExpr] String functions (length, contains, upper/lower, etc.).
  \item[AggExpr] Aggregations (sum, mean, count, min, max).
  \item[AliasExpr] Rename a result column.
\end{description}

\textbf{Type Promotion.} Binary operators check type compatibility and apply 
promotion rules:
\begin{itemize}
  \item Same type → same.
  \item Integer + integer → widest int (INT32 + INT64 → INT64).
  \item Integer + float → widest float (INT64 + FLOAT32 → FLOAT64).
  \item Numeric + string/bool → error.
\end{itemize}

\textbf{Eager Evaluation.} In eager mode, \texttt{EagerDataFrame::evalExpr} 
recursively evaluates the expression tree, dispatching each node onto 
Arrow Compute. For example, \texttt{col("x") + lit(3)} becomes 
\texttt{AddKernel(getColumn("x"), scalarFromLit(3))}.

\textbf{Lazy Evaluation.} In lazy mode, expressions are stored in the plan. 
At \texttt{collect()}, after optimization, the same \texttt{evalExpr} is 
called during execution.

\section{Logical Plan Representation}

The lazy plan is a DAG of \texttt{LogicalNode} subclasses:

\begin{description}
  \item[ScanNode] Leaf node. Reads a CSV or Parquet file. Carries optional 
    \texttt{projected\_columns} and \texttt{row\_limit} for optimization 
    annotations.
  \item[FilterNode] WHERE predicate.
  \item[SelectNode] SELECT expressions.
  \item[WithColNode] Add or replace one column.
  \item[GroupByNode] Declare grouping keys.
  \item[AggNode] Compute aggregations (paired with GroupByNode).
  \item[JoinNode] Equi-join. Carries \texttt{right}, \texttt{on}, \texttt{how}.
  \item[SortNode] ORDER BY columns.
  \item[LimitNode] LIMIT n rows.
  \item[SinkNode] Terminal write (CSV or Parquet).
\end{description}

Each node has an \texttt{execute()} method that takes eager input frames 
and returns an eager output frame. This is the hook that lets us reuse eager 
logic in the lazy executor.

\section{Query Optimizer}

The optimizer is a fixed-point rewriter: it applies five rules in sequence, 
checks structural plan equality (\texttt{planEquals}) against the previous 
iteration, and stops when the plan stabilizes or a safety cap (10 iterations) 
is hit.

\subsection{Optimization Rule: Predicate Pushdown}

\textbf{Purpose.} Move filters as close to their source as possible to reduce 
intermediate row counts.

\textbf{Transformation.} 

For a join: if a filter's predicate references only the left side,
\begin{verbatim}
Before: Filter(Join(L, R), pred)
After:  Join(Filter(L, pred), R)
\end{verbatim}

For a groupby: if a filter's predicate references only group keys,
\begin{verbatim}
Before: Filter(GroupBy(k, child), pred)
After:  GroupBy(k, Filter(child, pred))
\end{verbatim}

\textbf{Correctness.} For joins: rows on the left that fail \texttt{pred} 
cannot be part of any join output (predicate only depends on left), so dropping 
them below the join is safe. For groupby: within a group, all rows have the 
same key values, so a predicate on keys is constant within the group.

\textbf{Performance.} The join's hash table is built on fewer rows; groupby 
aggregates over fewer rows.

\subsection{Optimization Rule: Projection Pushdown}

\textbf{Purpose.} Only read columns actually used downstream.

\textbf{Transformation.} Top-down traversal maintains a required-set \(R\) of 
column names:

\begin{itemize}
  \item \texttt{SinkNode}: \(R\) unchanged.
  \item \texttt{SelectNode(exprs)}: \(R := \text{colRefs}(\text{exprs})\).
  \item \texttt{FilterNode(pred)}: \(R := R \cup \text{colRefs}(\text{pred})\).
  \item \texttt{JoinNode}: split \(R\) into left and right sides plus join keys.
  \item \texttt{ScanNode}: set \texttt{projected\_columns} = \(R \cap \text{file schema}\).
\end{itemize}

\textbf{Correctness.} We only ever drop columns never referenced upstream. 
Dropping unused columns cannot change the query result.

\textbf{Performance.} For a 12-column Parquet file, reading only 2 columns saves 
83\% of I/O bytes.

\subsection{Optimization Rule: Constant Folding}

\textbf{Purpose.} Evaluate fully-constant sub-expressions at plan time.

\textbf{Transformation.} If expression \(e\) contains no \texttt{ColExpr}, 
rewrite \(e \to \text{lit}(\text{evalExpr}(e))\).

\textbf{Example.}
\begin{verbatim}
Before: filter(col("age") > (lit(20) + lit(10)))
After:  filter(col("age") > lit(30))
\end{verbatim}

\textbf{Correctness.} Expression evaluation on a column-free tree produces 
the same Arrow scalar every time, so replacing it with that scalar is safe.

\textbf{Performance.} One addition per row eliminated; predicate is now a 
simple comparison.

\subsection{Optimization Rule: Expression Simplification}

\textbf{Purpose.} Apply algebraic identities to reduce expression complexity.

\textbf{Examples.}

\begin{tabular}{ll}
  \texttt{x + lit(0)} & $\to$ \texttt{x} \\
  \texttt{x * lit(1)} & $\to$ \texttt{x} \\
  \texttt{x \& lit(true)} & $\to$ \texttt{x} \\
  \texttt{\textasciitilde{}\textasciitilde{}x} & $\to$ \texttt{x} \\
  \texttt{x == x} & $\to$ \texttt{lit(true)} \\
\end{tabular}

\textbf{Correctness.} Each is a well-known identity. We specifically exclude 
null-unsafe rewrites: we do NOT rewrite \texttt{x - x → lit(0)} because 
\texttt{NULL - NULL = NULL}, not \texttt{0}.

\textbf{Performance.} Fewer kernel invocations; each identity removed is one 
compute call per row.

\subsection{Optimization Rule: Limit Pushdown}

\textbf{Purpose.} Move LIMIT below row-preserving operators and absorb it into 
scans.

\textbf{Transformation.}
\begin{verbatim}
LimitNode(n, SelectNode(...))  →  SelectNode(..., LimitNode(n, ...))
LimitNode(n, ScanNode(...))    →  ScanNode(..., row_limit=n)
\end{verbatim}

Does NOT push below FilterNode, JoinNode, SortNode, GroupByNode (these change 
row count or order).

\textbf{Correctness.} SelectNode and WithColNode are per-row transformations; 
they neither add nor remove rows. Reading only the first \(n\) rows of a scan 
produces exactly \(n\) rows.

\textbf{Performance.} On a 10M-row file, \texttt{head(5)} now reads only the 
first row group. On CSV, parsing stops after 5 lines.

\subsection{Optimizer Termination and Idempotence}

The optimizer is designed to terminate:
\begin{enumerate}
  \item Each rule is idempotent: once fully applied, it produces no change.
  \item The iteration count is capped at 10 (safety net).
  \item We check structural plan equality after each iteration; if unchanged, 
    we stop.
\end{enumerate}

\texttt{test\_main} includes a dedicated test (test 29) that verifies 
idempotence: running the optimizer twice on the same plan produces the same 
result.

\section{Execution}

\texttt{LazyDataFrame::collect()} does:
\begin{enumerate}
  \item Call \texttt{QueryOptimizer::optimize(plan\_)}.
  \item Call \texttt{executePlan(optimized\_root)} recursively:
    \begin{itemize}
      \item Materialize all child nodes.
      \item Call the node's \texttt{execute()} on eager input.
      \item Return the eager result.
    \end{itemize}
  \item Return the final \texttt{EagerDataFrame}.
\end{enumerate}

Each node's \texttt{execute()} delegates to the corresponding 
\texttt{EagerDataFrame} operator, ensuring we reuse all eager logic.

\section{Memory Management and Safety}

Every heap allocation is owned:
\begin{itemize}
  \item \texttt{std::shared\_ptr<T>} for reference-counted data (Arrow objects, 
    logical nodes, expressions).
  \item \texttt{std::unique\_ptr<T>} for exclusively-owned resources (Arrow 
    readers/writers).
  \item \texttt{std::string} / \texttt{std::vector} for byte/array data.
\end{itemize}

\texttt{EagerDataFrame} and \texttt{LazyDataFrame} are ``small handle'' types: 
one \texttt{shared\_ptr} member. Copies increment a refcount; moves are atomic 
pointer swaps.

\textbf{Exception safety.} Every Arrow API call is checked. Failures become 
\texttt{std::runtime\_error(status.ToString())}. Input validation throws 
\texttt{std::invalid\_argument} with descriptive messages. All state is held 
by smart pointers or standard containers, so unwinding is leak-free.

\textbf{ASAN + UBSan.} The full test suite runs under 
\texttt{-fsanitize=address,undefined} with \texttt{detect\_leaks=1}: zero 
leaks, zero UBSan errors.

\section{Testing}

The test suite (\texttt{tests/test\_main}) has 32 tests:

\begin{enumerate}
  \item Tests 1--10: Eager operators and schema handling.
  \item Tests 11--13: Expression evaluation, type system, null handling.
  \item Tests 14--18: Lazy execution, optimizer correctness.
  \item Tests 19--27: Edge cases (empty frames, 10k-row strings, parquet 
    round-trip, etc.).
  \item Tests 28--32: Optimizer idempotence, performance, and regression cases.
\end{enumerate}

\textbf{Benchmark.} Test 20 builds two 20k-row CSV files and a four-operator 
pipeline (scan → join → filter → select → head). It measures \texttt{collect\_raw()} 
(no optimizer) vs \texttt{collect()} (with optimizer). The minimum of 5 timed 
runs shows a typical speedup of \textbf{2×}.

\section{Code Organization}

\begin{description}
  \item[include/] Public API headers.
  \item[src/] Implementation.
  \item[src/internal/] Helper utilities (expression walking, Arrow bindings).
  \item[tests/] Test suite.
  \item[main.cpp] Demo program.
\end{description}

Total: approximately 2500 lines of C++, compact and readable.

\section{Conclusion}

DataFrameLib demonstrates a complete, working dataplane on top of Apache Arrow:
\begin{itemize}
  \item A Pandas-like eager API with zero-copy semantics.
  \item A Polars-like lazy API with a logical plan DAG.
  \item A fixed-point query optimizer with five industrial-strength rewrites.
  \item An executor that reuses eager logic for lazy queries.
  \item Comprehensive tests (32 tests, all passing).
  \item Full memory safety (ASAN/UBSan clean).
\end{itemize}

The codebase is intentionally small and readable, making it a good reference 
for teaching query optimization and demonstrating how modern data systems are 
structured.

\vspace{0.3in}
\noindent{\small Prepared for COP290 — Design Practices, Assignment 4}

\end{document}
5M-row table each removed identity saves one \texttt{Multiply} or
\texttt{Add} call.

\hypertarget{predicate-pushdown}{%
\subsubsection{5.3 --- Predicate Pushdown}\label{predicate-pushdown}}

\textbf{Description.} Move a \texttt{FilterNode} as close to its source
as possible. Two cases are handled: through a \texttt{JoinNode} (when
the predicate references only one side) and through a
\texttt{GroupByNode} (when the predicate only references grouping keys).

\textbf{Transformation.} If \texttt{pred} uses only columns of
\texttt{left} in \texttt{Filter(Join(left,\ right,\ on,\ how),\ pred)}:

\begin{verbatim}
Before:  Filter(Join(left, right), pred)
After:   Join(Filter(left, pred), right)
\end{verbatim}

(Analogous for right-side predicates on inner joins; outer joins only
allow pushdown on the preserved side.)

For \texttt{Filter(GroupBy(k,\ child),\ pred)} where
\texttt{colRefs(pred)\ ⊆\ k}:

\begin{verbatim}
Before:  Filter(GroupBy(k, child), pred)
After:   GroupBy(k, Filter(child, pred))
\end{verbatim}

\textbf{Correctness proof.} A row is kept by the outer filter iff
\texttt{pred} is true for the produced row. For an inner join on
\texttt{on}, a row on the left that fails \texttt{pred} cannot
participate in any produced row (because \texttt{pred} only depends on
left columns), so dropping it below the join changes nothing. For
\texttt{GroupBy} on keys \texttt{k}: rows in the same group share the
same values for \texttt{k}, so a predicate restricted to \texttt{k} is
constant within a group --- either all rows in the group pass or none
do, and moving the test above or below the grouping is equivalent.

\textbf{Concrete example.}

\begin{verbatim}
Before: scan(sales) ⨝ scan(stores)   → filter(store_size > 100)
After:  scan(sales) ⨝ (scan(stores) → filter(store_size > 100))
\end{verbatim}

\textbf{Performance benefit.} The right side of the join shrinks
\emph{before} the join builds its hash table, which for selective
filters can halve or quarter intermediate sizes.

\hypertarget{projection-pushdown}{%
\subsubsection{5.4 --- Projection Pushdown}\label{projection-pushdown}}

\textbf{Description.} Propagate the set of columns ultimately required
at the sink downward through the plan. At \texttt{ScanNode}, write the
intersection (with the file's real schema) into
\texttt{ScanNode::projected\_columns}, so the reader only materialises
those columns.

\textbf{Transformation.} Top-down traversal with a required-set
\texttt{R:\ optional\textless{}set\textless{}string\textgreater{}\textgreater{}}:

\begin{itemize}
\tightlist
\item
  \texttt{Sink\ /\ head\ /\ sort}: \texttt{R} unchanged.
\item
  \texttt{SelectNode(exprs)}: replace \texttt{R} with
  \texttt{colRefs(exprs)}.
\item
  \texttt{FilterNode(pred)}: \texttt{R\ :=\ R\ ∪\ colRefs(pred)}.
\item
  \texttt{WithColNode(name,expr)}:
  \texttt{R\ :=\ (R\ \textbackslash{}\ \{name\})\ ∪\ colRefs(expr)}.
\item
  \texttt{GroupByNode\ +\ AggNode}:
  \texttt{R\ :=\ keys\ ∪\ colRefs(aggregations)}.
\item
  \texttt{JoinNode(on,\ left,\ right)}: split \texttt{R} into
  \texttt{R\_left\ ∪\ on} and \texttt{R\_right\ ∪\ on}, descend into
  each side.
\item
  \texttt{ScanNode}: record
  \texttt{projected\_columns\ =\ R\ ∩\ file\_schema}.
\end{itemize}

\textbf{Correctness proof.} Every node in the plan outputs a subset of
some set of columns. The transformation only ever \emph{drops} columns
that are never referenced upstream of the scan; dropping unused columns
cannot change the result of a relational operator.

\textbf{Concrete example.}

\begin{verbatim}
sales.csv has 12 columns. Query only needs (id, amount).
Before: scan(sales) → select(id, amount)
After:  scan(sales, projected={id, amount}) → select(id, amount)
\end{verbatim}

\textbf{Performance benefit.} For a 12-column Parquet file this reads
2/12 ≈ 17\% of the bytes. Even for CSV, converting 10 unused columns per
row to Arrow arrays is skipped.

\hypertarget{limit-pushdown}{%
\subsubsection{5.5 --- Limit Pushdown}\label{limit-pushdown}}

\textbf{Description.} Move \texttt{LimitNode(n)} below every operator
that preserves row count and order. In particular, \texttt{LimitNode} is
swapped below \texttt{SelectNode} and \texttt{WithColNode}, and absorbed
into \texttt{ScanNode::row\_limit} when it reaches a scan.

\textbf{Transformation.}

\begin{verbatim}
LimitNode(n, SelectNode(s,  child))  →  SelectNode(s,  LimitNode(n, child))
LimitNode(n, WithColNode(w, child))  →  WithColNode(w, LimitNode(n, child))
LimitNode(n, ScanNode(path,...))     →  ScanNode(path, ..., row_limit=n)
\end{verbatim}

Does \emph{not} push below \texttt{FilterNode}, \texttt{JoinNode},
\texttt{SortNode}, \texttt{GroupByNode}, \texttt{AggNode} (these change
row count / order).

\textbf{Correctness proof.} \texttt{SelectNode} and \texttt{WithColNode}
are per-row transformations --- they neither add, remove, nor reorder
rows. \texttt{select(head(x,\ n))\ =\ head(select(x),\ n)} because both
evaluate the select on the same set of n rows in the same order. For
\texttt{ScanNode}, reading only the first \texttt{n} rows of the source
file is exactly what the unoptimised \texttt{LimitNode} would produce.

\textbf{Concrete example.}

\begin{verbatim}
Before: scan(big.csv) → with_column(c, a+b) → head(5)
After:  scan(big.csv, row_limit=5) → with_column(c, a+b)
\end{verbatim}

\textbf{Performance benefit.} On a 10M-row Parquet file taking
\texttt{head(5)} now reads only the first row group. In CSV mode we stop
reading after \texttt{n} lines. The \texttt{with\_column} compute also
runs on 5 rows instead of 10M.

\hypertarget{memory-management}{%
\subsection{6. Memory Management}\label{memory-management}}

Every heap allocation in the library is owned by a smart pointer:

\begin{itemize}
\tightlist
\item
  \texttt{std::shared\_ptr\textless{}T\textgreater{}} for values that
  may have multiple owners: \texttt{arrow::Table},
  \texttt{arrow::ChunkedArray}, \texttt{LogicalNode}, \texttt{Expr}.
\item
  \texttt{std::unique\_ptr\textless{}T\textgreater{}} for
  exclusively-owned Arrow readers / writers inside the I/O code.
\item
  \texttt{std::string} / \texttt{std::vector} everywhere for owning byte
  / array data --- no \texttt{char*} escape hatches.
\end{itemize}

\texttt{EagerDataFrame} and \texttt{LazyDataFrame} are both ``small
handle'' types: the only data member is one \texttt{shared\_ptr}, so
copies are \texttt{atomic\_fetch\_add(1)} on a refcount and moves are a
pointer swap. Both classes have the full defaulted rule-of-5
(\texttt{=\ default}) plus \texttt{noexcept} move operators.

Arrow file handles are never manually closed. Readers returned by
\texttt{arrow::io::ReadableFile::Open} are owned by reader objects whose
destructors close them --- RAII all the way.

The only filesystem resource we create directly is the intermediate
\texttt{.dot} file in \texttt{explain()}. This is wrapped in a local
\texttt{TempDotFile} struct whose destructor deletes the file unless a
\texttt{keep} flag is set (e.g.~when \texttt{dot} is missing and we want
the user to see the path). \texttt{TempDotFile} is non-copyable to
prevent double-delete.

\textbf{Exception safety.} Every Arrow API call that returns
\texttt{arrow::Status} /
\texttt{arrow::Result\textless{}T\textgreater{}} is checked; failures
are converted to \texttt{std::runtime\_error(status.ToString())}. Input
validation throws \texttt{std::invalid\_argument} with a descriptive
message (missing column name, unknown join mode, non-boolean predicate,
\ldots). Because all owning state is held by smart pointers and standard
containers, unwinding is leak-free.

The test suite is run under \texttt{-fsanitize=address,undefined} with
\texttt{detect\_leaks=1}; the complete 31-test run reports zero leaks
and zero UBSan errors.

\hypertarget{testing}{%
\subsection{7. Testing}\label{testing}}

The test binary \texttt{tests/test\_main} contains 31 tests divided into
five categories:

\begin{enumerate}
\def\labelenumi{\arabic{enumi}.}
\tightlist
\item
  \textbf{Schema \& eager operators} (tests 1--10): \texttt{read\_csv},
  \texttt{from\_columns}, \texttt{select}, \texttt{filter},
  \texttt{with\_column}, \texttt{group\_by\ +\ \ \ \ aggregate},
  \texttt{join}, \texttt{sort}, \texttt{head}, \texttt{write\_csv}.
\item
  \textbf{Expression \& type system} (tests 11--13): arithmetic
  promotion, null handling, string functions, incompatible-type
  detection.
\item
  \textbf{Lazy execution} (tests 14--15): plan construction,
  \texttt{collect()}, \texttt{explain()}, parity between eager and lazy
  pipelines.
\item
  \textbf{Optimiser} (tests 16--20): one test per rule plus a benchmark.
  Each optimiser test asserts both \emph{structural} correctness
  (inspect the optimised plan) and \emph{semantic} correctness
  (\texttt{collect()\ ==\ collect\_raw()}).
\item
  \textbf{Edge cases} (tests 21--31): empty frames, single-row
  aggregations, all-null columns, 10 000-row string columns, chained
  filters, multi-key joins, 8-operation lazy plans, parquet round-trip,
  optimiser idempotency, out-of-range \texttt{head}, and
  \texttt{std::invalid\_argument} on string+int arithmetic.
\end{enumerate}

\hypertarget{benchmark-methodology}{%
\subsubsection{Benchmark methodology}\label{benchmark-methodology}}

The optimiser benchmark (test 20) builds two 20 000-row CSV files and a
four-operator pipeline
(\texttt{scan\ →\ join\ →\ filter\ →\ select\ →\ head}). It measures
\texttt{collect\_raw()} (no optimiser) vs \texttt{collect()} (optimiser
on) using a \texttt{minTimeMs} helper that runs 3 warm-up iterations
followed by 5 timed iterations and reports the minimum --- the minimum
is a much more stable metric than the mean for short measurements. The
test asserts the optimised run is at least 10 \% faster; in practice the
speedup is \textasciitilde2×.

Memory safety is verified by recompiling the suite with
\texttt{-fsanitize=address,undefined} and
\texttt{ASAN\_OPTIONS=detect\_leaks=1}, which reports no errors and no
leaks across all 31 tests.

\end{document}
